\subsection*{First Example: Tiled Addition Through shared storage}

\begin{lstlisting}[language=croqtile,caption={First Example: Tiled Addition Through shared storage}]
__co__ s32 [64, 128] tiled_add_2d(s32 [64, 128] lhs, s32 [64, 128] rhs) {
  s32 [lhs.span] output;
  parallel {tr, tc} by [4, 8] {
    lhs_load = dma.copy lhs.chunkat(tr, tc) => shared;
    rhs_load = dma.copy rhs.chunkat(tr, tc) => shared;
    foreach {i, j} in lhs_load.span
      output.at(tr # i, tc # j) =
          lhs_load.data.at(i, j) + rhs_load.data.at(i, j);
  }
  return output;
}
\end{lstlisting}

The arithmetic is still element-wise addition. What changed is the data path:

\begin{enumerate}[leftmargin=1.6em]
\item \texttt{chunkat(tr, tc)} selects one logical tile from each input.
\item \texttt{dma.\allowbreak{}copy .\allowbreak{}.\allowbreak{}.\allowbreak{} =\textgreater{} shared} stages that tile into block-visible fast memory.
\item \texttt{tr \# i} and \texttt{tc \# j} map tile-local coordinates back to global output coordinates.
\end{enumerate}

This is the core Croqtile pattern: select a tile, move it, compute on the moved tile.

\begin{figure}[htbp]\centering
\includegraphics[width=0.86\linewidth,height=0.42\textheight,keepaspectratio]{f04-fig2_tiled_add_light.png}
\end{figure}
\begin{center}\small\emph{Load a tile, compute on the staged data, then write results back using global coordinates.}\end{center}
